%%%%%
%@TheDoctorRAB
%%%%%
%
%%%%%
%standard report format for journal papers
%%%%%
%
%%%%%
%PDF tagging accessible file
%compile with pdflatex
%run $show-pdf-tags --xml filename.pdf
%%%%%
%
%%%%%
%REFERENCES
%neup.bst - numbered citations in order of appearance, short author list with et al in reference section
%nsf.bst - numbered citations in order of appearance, full author list in references section
%standard.bst - citations with author last name with et al for more than 2 authors; full author list in references section
%ans.bst is for ANS only. 
%
%author = {Lastname, Firstname and Lastname, Firstname and Lastname, Firstname} for all bst formats
%bst renders the author list itself
%
%author = {{Nuclear Regulatory Commission}} if the author is an organization, institution, etc., and not people
%
%title = {{}} for all
%
%for all - use \citep{-} - [1] or (Borrelli, 2021) in the text
%standard.bst \cite{-} - Borrelli (2021) in the text
%standard.bst lists references alphabetically
%the rest list numerically
%%%%%


%%%%% enable tagging for accessibility
\RequirePackage{pdfmanagement-testphase}
\DocumentMetadata{
    tagging=on,
    pdfversion=2.0,
    pdfstandard=ua-2,
    lang=en-US,
    debug={log=all},
    testphase={latest,tagpdf},
    tagging-setup={math/setup={mathml-AF,mathml-SE}}
}
\tagpdfsetup{math/alt/use}
%%%%%


%%%%% general 
\documentclass[11pt,a4paper]{article}
\usepackage[lmargin=1in,rmargin=1in,tmargin=1in,bmargin=1in]{geometry}
\usepackage[pagewise]{lineno} %linenumbering restarts each page
%\usepackage[running]{lineno} %linenumbering continuous
\usepackage{mmap} %helps make pdfs copy and paste
\usepackage{setspace}
\usepackage{ulem} %strikethrough - do not \sout{\cite{}}
\usepackage[pdftex,dvipsnames]{xcolor,colortbl} %change font color
\usepackage{graphicx}
\usepackage{tablefootnote}
\usepackage{footnotehyper}
\usepackage{float}
%\usepackage{mypythonhighlight,verbatim}
%\usepackage{subfig}
\usepackage[yyyymmdd]{datetime} %date format
\renewcommand{\dateseparator}{.}
%\graphicspath{{$TEXIMG/}} %path to graphics
\setcounter{secnumdepth}{5} %set subsection to nth level
\usepackage{needspace}
\usepackage[stable,hang,flushmargin]{footmisc} %footnotes in section titles and no indent; standard.bst
\usepackage[inline]{enumitem} %https://www.tug.org/texlive//Contents/live/texmf-dist/doc/latex-dev/latex-lab/latex-lab-enumitem.pdf
\usepackage{boldline}
\usepackage{makecell}
\usepackage{booktabs}
\usepackage{amssymb}
\usepackage{gensymb}
\usepackage{amsmath,nicefrac}
\usepackage{physics}
\usepackage{lscape}
\usepackage{array}
\usepackage{chngcntr}
\usepackage{sectsty}
\usepackage{textcomp}
\usepackage{lastpage}
\usepackage{xargs} %for \newcommandx
\usepackage[colorinlistoftodos,prependcaption,textsize=tiny]{todonotes} %makes colored boxes for commenting
\usepackage{soul}
\usepackage{color}
\usepackage{marginnote}
\usepackage[figure,table]{totalcount}
\usepackage{parskip}
%%%%%


%%%%% links and tags
\usepackage{hyperref}
\hypersetup{
    colorlinks,
    linkcolor=black,
    citecolor=black,
    urlcolor=blue,
    pdftitle={Homework},
    pdfauthor={R. A. Borrelli},
    pdfkeywords={homework},
    pdflanguage={en-US},
    pdfsubject={homework},
    pdfdisplaydoctitle=true 
} 
\usepackage[capitalise]{cleveref} %must be loaded after hyperref
%%%%%


%%%%% tikz
\usepackage{pgf}
\usepackage{tikz} % required for drawing custom shapes
\usetikzlibrary{shapes,arrows,automata,trees}
%%%%%


%%%%% fonts
\usepackage{times}
%arial - uncomment next two lines
%\usepackage{helvet}
%\renewcommand{\familydefault}{\sfdefault}
%%%%%


%%%%% references
%\usepackage[round,semicolon]{natbib} %for (Borrelli 2021; Clooney 2019) - standard.bst 
\usepackage[numbers,sort&compress]{natbib} %for [1-3] - nsf.bst, neup.bst
\setlength{\bibsep}{7pt} %sets space between references
%\renewcommand{\bibsection}{} %suppresses large 'references' heading
%\renewcommand\bibpreamble{\vspace{\baselineskip}} %sets spacing after heading if not using default references heading
%%%%%


%%%%% tables and figures
\usepackage{longtable} %need to put label at top under caption then \\ - use spacing
\usepackage{tablefootnote}
\usepackage{tabularx}
\usepackage{multirow}
\usepackage{tabto} %general tabbed spacing
\usepackage{pdfpages}
\usepackage{wrapfig} %wraps figures around text
\setlength{\intextsep}{0.00mm}
\setlength{\columnsep}{1.00mm}
\usepackage[singlelinecheck=false,labelfont=bf]{caption}
\usepackage{subcaption}
\captionsetup[table]{justification=centering,skip=5pt,labelformat={default},labelsep=period,name={Table}} %sets a space after table caption
\captionsetup[figure]{justification=justified,skip=5pt,labelformat={default},labelsep=period,name={Figure}} %sets space above caption, 'figure' format
\captionsetup[wrapfigure]{justification=centering,aboveskip=0pt,belowskip=0pt,labelformat={default},labelsep=period,name={Fig.}} %sets space above caption, 'figure' format
\captionsetup[wraptable]{justification=centering,aboveskip=0pt,belowskip=0pt,labelformat={default},labelsep=period,name={Table}} %sets space above caption, 'figure' format
%%%%%


%%%%% watermark
%\usepackage[firstpage,vpos=0.63\paperheight]{draftwatermark}
%\SetWatermarkText{\shortstack{DRAFT\\do not distribute}}
%\SetWatermarkScale{0.20}
%%%%%


%%%%% cross referencing files
%\usepackage{xr} %for revisions - will cross reference from one file to here
%\externaldocument{/path/to/auxfilename} %aux file needed
%%%%%


%%%%% toc and glossaries
\usepackage[toc,title]{appendix}
\usepackage[acronym,nomain,nonumberlist]{glossaries}
\makenoidxglossaries
%do not load titlesec,titletoc
%%%%%


%%%%% editing
\newcommand{\edit}[1]{\textcolor{blue}{#1}} %shortcut for changing font color on revised text
\newcommand{\fn}[1]{\footnote{#1}} %shortcut for footnote tag
\newcommand*\sq{\mathbin{\vcenter{\hbox{\rule{.3ex}{.3ex}}}}} %makes a small square as a separator $\sq$
%\newcommand{\sk}[1]{\sout{#1}} %shortcut for default strikethrough - do not sk through citep
\newcommand\sk{\bgroup\markoverwith{\textcolor{red}{\rule[0.5ex]{1pt}{1pt}}}\ULon} %strikethrough with red line; not in \section{}
%\st{} does strikethrough using soul package but does not like acronyms
\newcommand{\blucell}{\cellcolor{aliceblue}} %use to shade in table cell
\newcommand{\grycekk}{\cellcolor{lightgray}} %use to shade in table cell
\newcommand{\whicell}{\cellcolor{antiquewhite}} %use to shade in table cell
%%%%%


%%%%% colors
%http://latexcolor.com/
%https://en.wikibooks.org/wiki/LaTeX/Colors#:~:text=black%2C%20blue%2C%20brown%2C%20cyan,be%20available%20on%20all%20systems.
\definecolor{aliceblue}{rgb}{0.94, 0.97, 1.0}
\definecolor{antiquewhite}{rgb}{0.98, 0.92, 0.84}
\definecolor{lightmauve}{rgb}{0.86, 0.82, 1.0}
\definecolor{brilliantlavender}{rgb}{0.96, 0.73, 1.0}
\definecolor{brandeisblue}{rgb}{0.0, 0.44, 1.0}
\definecolor{darkmidnightblue}{rgb}{0.0, 0.2, 0.4}

\newcommand{\x}{\cellcolor{aliceblue}} %use to shade in table cell
\newcommand{\y}{\cellcolor{lightgray}} %use to shade in table cell
\newcommand{\z}{\cellcolor{antiquewhite}} %use to shade in table cell
%%%%%


%%%%% acronyms
\newcommand{\acf}{\acrfull} %full acronym
\newcommand{\acl}{\acrlong} %long acronym
\newcommand{\acs}{\acrshort} %short acronym

\newcommand{\acfp}{\acrfullpl} %full acronym plural
\newcommand{\aclp}{\acrlongpl} %long acronym plural
\newcommand{\acsp}{\acrshortpl} %short acronym plural
%%%%%


%%%%% todonotes
\newcommandx{\cmt}[2][1=]{\todo[author=\textbf{STRUCTURE},tickmarkheight=0.15cm,linecolor=red,backgroundcolor=red!25,bordercolor=black,#1]{#2}}
\newcommandx{\con}[2][1=]{\todo[author=\textbf{CONTENT},tickmarkheight=0.15cm,linecolor=brilliantlavender,backgroundcolor=brilliantlavender,bordercolor=black,#1]{#2}}
\newcommandx{\rab}[2][1=]{\todo[noline,author=\textbf{RAB},backgroundcolor=Plum!25,bordercolor=black,#1]{#2}}


%\newcommandx{\jon}[2][1=]{\todo[noline,author=\textbf{ATTN: Johnson},backgroundcolor=blue!25,bordercolor=black,#1]{#2}}
%\newcommandx{\han}[2][1=]{\todo[noline,author=\textbf{ATTN: Haney},backgroundcolor=OliveGreen!25,bordercolor=black,#1]{#2}}
%\newcommandx{\rab}[2][1=]{\todo[author=\textbf{RAB},tickmarkheight=0.15cm,linecolor=Plum,backgroundcolor=Plum!25,bordercolor=black,#1]{#2}}
%\newcommandx{\han}[2][1=]{\todo[author=\textbf{ATTN: Haney},tickmarkheight=0.15cm,linecolor=OliveGreen,backgroundcolor=OliveGreen!25,bordercolor=OliveGreen,#1]{#2}}
%\newcommandx{\jon}[2][1=]{\todo[author=\textbf{ATTN: Johnson},tickmarkheight=0.15cm,linecolor=blue,backgroundcolor=blue!25,bordercolor=blue,#1]{#2}}


% highlighting 
\DeclareRobustCommand{\hlc}[1]{{\sethlcolor{LimeGreen}\hl{#1}}}
\makeatletter
    \if@todonotes@disabled
    \newcommand{\hlh}[2]{#1}
    \else
    \newcommand{\hlh}[2]{\han{#2}\hlc{#1}}
    \fi
    \makeatother

\DeclareRobustCommand{\hld}[1]{{\sethlcolor{CornflowerBlue}\hl{#1}}}
\makeatletter
    \if@todonotes@disabled
    \newcommand{\hlj}[2]{#1}
    \else
    \newcommand{\hlj}[2]{\jon{#2}\hld{#1}}
    \fi
    \makeatother

\DeclareRobustCommand{\hlf}[1]{{\sethlcolor{lightmauve}\hl{#1}}}
\makeatletter
    \if@todonotes@disabled
    \newcommand{\hlb}[2]{#1}
    \else
    \newcommand{\hlb}[2]{\rab{#2}\hlf{#1}}
    \fi
    \makeatother
%%%%%


%%%%% table alignments
\newcolumntype{L}[1]{>{\raggedright\let\newline\\\arraybackslash\hspace{0pt}}m{#1}} %uses \raggedright with m,p{} in table column
\newcolumntype{C}[1]{>{\centering\let\newline\\\arraybackslash\hspace{0pt}}m{#1}} %uses \raggedright with m,p{} in table column
\newcolumntype{R}[1]{>{\raggedleft\let\newline\\\arraybackslash\hspace{0pt}}m{#1}} %uses \raggedright with m,p{} in table column
%%%%%


%%%%% table contents
\makeatletter
\renewcommand\tableofcontents{%
    \@starttoc{toc}%
}
\makeatother

\makeatletter
\renewcommand\listoffigures{%
    \@starttoc{lof}%
}
\makeatother

\makeatletter
\renewcommand\listoftables{%
    \@starttoc{lot}%
}
\makeatother

\makeatletter
\newcommand*\ftp{\fontsize{16.5}{17.5}\selectfont}
\makeatother
%%%%%


%%%%% header and footer
\usepackage{fancyhdr}
\pagestyle{fancy}
\fancyhf{} %move page number to bottom right
%\renewcommand{\headrulewidth}{0pt} %turn off line in header
\lhead{\scriptsize NE585}
\chead{\scriptsize \textit{Project 6 -- Back end of the nuclear fuel cycle}}
\rhead{\scriptsize \today}
\rfoot{\thepage}
%%%%%


%%%%% acronyms
\input{$ACRONYM/acronyms}
%\input{../acronyms/acronyms}
%%%%%


\begin{document}

\begin{titlepage}
    \title{
        NE585 -- Nuclear fuel cycle analysis\\
        Project 6 -- Back end of the nuclear fuel cycle\\
    }
    \author{
        Name
        \\ \\ \\
        University of Idaho $\sq$ Idaho Falls Center for Higher Education
        \\ \\
        Nuclear Engineering and Industrial Management Department
        \\ \\ \\
        email 
    }
\clearpage %not have page number on title page
\maketitle
\vspace*{\fill}
\begin{flushright}{
        Total -- 275
}
\end{flushright}
\thispagestyle{empty} %start with page number 1 on second page
\end{titlepage}


\section{Finding another repository -- (25)}
The United States has about 70000 metric tons of used fuel, stored onsite at the reactors. Even if Yucca Mountain opened tomorrow, it would be filled immediately and there would still be used fuel left over. Therefore, a second and third repository are going to be needed. Reprocessing aside, which significantly reduces waste volume, and given the high-level geologic criteria for a repository, where might a second or third repository be sited in the United States and why? Given the culture of the local population, how likely would they consent to a repository? Would your consent-based approach from PROJECT 3 be useful to this end?





\newpage

\section{Geologic analogue -- (20)}
Is there a natural analogue for a geologic repository? 





\newpage

\section{Linear programming -- (25)}
Use the linear programming procedure to maximize the waste oxide mass for the glass waste form of the Japanese concept, subject to the six constraints discussed in class.





\newpage

\section{Rokkasho capacity -- (25)}
78 kg of waste oxide is produced per metric ton of used fuel (MTU) during reprocessing. If there are 15,000 MTU in Japan right now, how many waste canisters are needed (based on the answer obtained previously)? The Rokkasho reprocessing plant can hold 2880 canisters at any one time. Would the plant have enough space to store the canisters? If not, what should be done? 





\newpage

\section{Rokkasho process modeling -- (30)}
Rokkasho can store 20,000 MTU on the site and process 800 MTU/year. A typical reactor produces 20 MTU/year, and Japan has 50 reactors. 1500 MTU can be stored onsite at each reactor. Derive a process model for the production of canisters per year. Assume the current 15000 MTU is stored equally at each plant at TIME = 0, so Rokkasho starts off totally empty of used fuel and canisters. Make any other reasonable assumptions. Produce several useful plots from the model simulation. Based on the results of the model, offer some insight into a waste management plan/policy for Japan. 





\newpage

\section{Into eternity -- (50)}
Write an op-ed style piece (750 - 1000 words) on `Into Eternity.' Do you think the safety of a repository is feasible centuries into the future? What would you say to future society about the repository? Write a short op-ed piece (500 words or so) on ‘Into Eternity.’ Do you think the safety of a repository is feasible centuries into the future? What would you say to future society about the repository?


\textit{We will verify if the movie is available first.}





\newpage

\section{Consent-based approach -- (50)}
Everyone agrees that a `consent-based approach' must be used when siting any nuclear facility. However, no one really says what this would look like. From an engineering perspective, we want models and easily quantifiable metrics. Within this context, what might a consent-based approach look like?





\newpage

\section{Yucca mountain -- (50)}
Give an update on the current status of the Yucca mountain geologic repository. Will it ever become a repository? What else could the facility/property be used for?





%%%%% header and footer
\pagestyle{fancyplain}
\fancyhf{} %clear header and footer 
\renewcommand{\headrulewidth}{0pt} %set line thickness in header; uncomment as is to remove line
\rfoot{\fancyplain{}{\thepage}}
\nolinenumbers
%%%%%


%%%%%%% citations
%
%%% make sure $BIBREF contains the path to each bib file in .bashrc
%%% make sure Overleaf has the latexmkrc pointing to where the bib and bst are located
%
\bibliographystyle{nsf}
\setlength{\bibhang}{0pt}
\bibliography{
}
%
%%%%%%%


\newpage


\section*{Tables}

{%
\let\oldnumberline\numberline%
\renewcommand{\numberline}{\tablename~\oldnumberline}%
\listoftables%
}


\newpage


%\tagpdfsetup{table/header-rows=1}
%\begin{longtable}{|c|c|}
%        \caption{Average counts for unknown radioactive sample}
%        \label{tab-half-life}
%        \\
%        \hline
%        \textbf{Time (s)}& 
%        \textbf{Counts per second (cps)}
%        \\
%        \hline
%        \endfirsthead
%        0
%        &1000
%        \\
%        30
%        &856.7
%        \\
%        60
%        &734.0
%        \\
%        90
%        &628.8
%        \\
%        120
%        &538.8
%        \\
%        150
%        &461.6
%        \\
%        180
%        &395.4
%        \\
%        210
%        &338.8
%        \\
%        240
%        &290.3
%        \\
%        270
%        &248.7
%        \\
%        300
%        &213.0
%        \\
%        330
%        &182.5
%        \\
%        360
%        &156.4
%        \\
%        390
%        &134.0
%        \\
%        420
%        &114.8
%        \\
%        450
%        &98.3
%        \\
%        480
%        &84.2
%        \\
%        510
%        &72.2
%        \\
%        540
%        &61.8
%        \\
%        570
%        &53.0
%        \\
%        600
%        &45.4
%        \\
%        \hline
%\end{longtable}


\newpage


\section*{Figures}

{%
\let\oldnumberline\numberline%
\renewcommand{\numberline}{\figurename~\oldnumberline}%
\listoffigures%
}


\newpage


%\begin{figure}[h!]
%    \centering
%    \includegraphics[width=0.X\textwidth]{filename}
%    \caption{Descriptive caption}
%    \label{fig-label-name}
%\end{figure}


\begin{appendices}


%%%%% formatting
    \renewcommand{\thesection}{\Roman{section}}
%%%%%


%    \section{\acs{mcp} input decks} \label{app-mcnp}


\newpage


    \section{Acronyms}
    \printnoidxglossary[title={}]


\end{appendices}


\end{document}
